%% 第 16 章：伴随算子
%% 本文件由 tools/md2tex.py 自动生成，请勿直接编辑。

\chapter{伴随算子}


\section{为什么需要定义伴随}

设：

\[
T:V\to W.
\]

若固定 $w\in W$，则：

\[
v\mapsto\langle T(v),w\rangle_W
\]

是 $V$ 上的线性泛函。

它表示：先将 $v$ 通过 $T$ 送到 $W$，再用 $w$ 所代表的内积测量方式对其测量。

\section{伴随的定义}

对固定的 $w\in W$，考虑：

\[
v\mapsto\langle T(v),w\rangle_W.
\]

由于 $T$ 是线性映射，且内积对第一个位置线性，这确实是一个定义在 $V$ 上的线性泛函。

由 Riesz 表示定理，$V$ 上的每个线性泛函都能被唯一写成“与 $V$ 中某个向量做内积”的形式。因此，存在唯一的向量 $T^*(w)\in V$，使：

\[
\langle T(v),w\rangle_W
=
\langle v,T^*(w)\rangle_V.
\]

这里，$T^*(w)$ 就是将 $W$ 中由 $w$ 给出的测量方式，经由 $T$ 拉回到 $V$ 后所对应的唯一向量。

伴随与此前的对偶映射本质上都是对于泛函这种测量方式在输入输出空间中的转换。

但区别在于：

\begin{jitem}
\item %
对偶映射将 $W$ 中的泛函拉回为 $V$ 中的泛函；也就是说，$T'(\varphi)=\varphi\circ T.$
\item %
伴随利用 Riesz 表示定理，将这个被拉回去的泛函，直接用向量而不是复合的映射的方式，在 $V$ 中表示出来，并且表示为一种内积的形式。
\end{jitem}

\begin{jdefn}{伴随算子的定义}{r:17:1}
伴随算子是唯一满足下式的线性映射：

\[
T^*:W\to V,
\]

\[
\langle T(v),w\rangle_W
=
\langle v,T^*(w)\rangle_V.
\]
\end{jdefn}

\begin{jprop}{伴随算子的存在与唯一性}{r:17:2}
对于每个 $w\in W$，映射：

\[
v\mapsto\langle T(v),w\rangle_W
\]

是 $V$ 上的线性泛函。

Riesz 表示定理保证存在唯一向量 $T^*(w)$ 表示这个泛函，因此伴随存在且唯一。

$T^*$ 具有加性和齐性，因此 $T^*$ 是线性映射。
\end{jprop}

\begin{jprop}{伴随的代数性质}{r:17:3}
设：

\[
S,T:V\to W,
\]

且 $\lambda\in\mathbb F$。

则：

\[
(S+T)^*=S^*+T^*,
\]

\[
(\lambda T)^*=\overline{\lambda}T^*,
\]

\[
(T^*)^*=T.
\]

若：

\[
T:U\to V,
\qquad
S:V\to W,
\]

则：

\[
(S\circ T)^*=T^*\circ S^*.
\]
\end{jprop}

\begin{jproof}{这些性质的证明思路}
证明使用同一种思路：将等式两边放进内积，并不断使用伴随定义把算子从内积的第一个位置移到第二个位置。

加法对应内积的加法性。

标量取共轭，是因为本笔记采用第一位置线性的内积约定：标量从第二个位置移出时会变为共轭。

复合反序只需要通过反复应用伴随的定义在内积的表达式上即可。
\end{jproof}
